% swim-times-defect-report.tex
% Defect investigation and corrective-action report for swim-times.
% Compiled with: pdflatex swim-times-defect-report.tex

\documentclass[11pt]{article}

\usepackage{geometry}
\geometry{letterpaper, margin=1in}

\usepackage[T1]{fontenc}
\usepackage{hyperref}
\usepackage{booktabs}
\usepackage{array}
\usepackage{longtable}
\usepackage{enumitem}
\usepackage{titlesec}
\usepackage{fancyhdr}
\usepackage{xcolor}
\usepackage{listings}

\definecolor{castblue}{RGB}{0,48,135}
\definecolor{shade}{RGB}{245,246,248}
\definecolor{bad}{RGB}{150,30,30}
\definecolor{good}{RGB}{20,100,50}

\hypersetup{colorlinks=true, linkcolor=castblue, urlcolor=castblue,
            citecolor=castblue}

\pagestyle{fancy}
\fancyhf{}
\renewcommand{\headrulewidth}{0.4pt}
\fancyhead[L]{\small College Area Swim Team}
\fancyhead[R]{\small Defect Report: swim-times}
\fancyfoot[C]{\thepage}

\titleformat{\section}{\large\bfseries\color{castblue}}{\thesection}{0.5em}{}[\titlerule]
\titleformat{\subsection}{\normalsize\bfseries}{\thesubsection}{0.5em}{}

\newcommand{\code}[1]{\texttt{#1}}
\newcommand{\req}[1]{\texttt{#1}}

\lstset{
  basicstyle=\ttfamily\small,
  backgroundcolor=\color{shade},
  frame=single,
  rulecolor=\color{gray!40},
  framesep=6pt,
  breaklines=true,
  columns=fullflexible,
  showstringspaces=false,
  xleftmargin=0pt,
  literate={-}{-}1
}

\begin{document}

% -----------------------------------------------------------------
\begin{titlepage}
  \centering
  \vspace*{2cm}
  {\Huge\bfseries\color{castblue} Defect Investigation and\\[0.3em]
   Corrective Action Report\\[0.8em]}
  {\Large \code{swim-times}: ``the program appears to be not working''}\\[2em]
  {\large College Area Swim Team (CAST)}\\[0.5em]
  {\normalsize Prepared for: CAST Coaching Staff and Maintainers}\\[0.5em]
  {\normalsize Prepared by: CAST Technology Working Group}\\[2em]
  \begin{tabular}{ll}
    Document identifier: & CAST-DR-001                                   \\
    Revision:            & 5.0                                           \\
    Date:                & September 18, 2026                            \\
    Status:              & Released                                      \\
    Originating request: & \code{requirements9.txt} item 1;              \\
                         & \code{requirements10.txt} items 1--3;         \\
                         & ``use the database on 100.77.243.69 only'';   \\
                         & ``how do I add new swimmers'';                \\
                         & ``analyze addswimuser; make swim-times work   \\
                         & with new users and aliases''                  \\
    Subject build:       & \code{swim-times.w} rev.\ 6.0                 \\
  \end{tabular}
  \vfill
  {\small\textit{Prepared in accordance with the principles of rational
  documentation articulated in D.\,L.\,Parnas and P.\,C.\,Clements,
  ``A Rational Design Process: How and Why to Fake It,''
  \textit{IEEE Transactions on Software Engineering}, vol.\,SE-12,
  no.\,2, pp.\,251--257, February 1986.  In the spirit of that paper the
  investigation is presented in the order a reader needs it, not in the
  order it actually happened.}}
\end{titlepage}

\tableofcontents
\clearpage

% =================================================================
\section{Executive summary}
% =================================================================

\code{swim-times} stopped returning data.  Every invocation that
touched the network ended in a single uninformative line:

\begin{lstlisting}
$ ./swim-times -o stella
Error: person lookup request failed
\end{lstlisting}

\noindent The investigation found \textbf{three distinct problems},
layered one on top of another.  Only the first is an external change;
the other two are defects in \code{swim-times} itself, and the second
of them is what made the first so hard to see.

\begin{enumerate}[leftmargin=*]
\item \textbf{The USA Swimming Data Hub withdrew anonymous access to
      its times endpoints.}  Member search and every times query now
      answer \code{403 Forbidden} to a caller who is not signed in.
      This is a policy change at the service and cannot be repaired
      from this side; it can only be adapted to.
\item \textbf{\code{http\_request} never looked at the HTTP status
      code.}  It reported failure by returning \code{NULL}, and it
      returned \code{NULL} for any response with an empty body.  A
      \code{403} with \code{content-length:\,0} was therefore
      indistinguishable from a dead network.  The one piece of
      information that would have explained the outage was discarded
      unread.
\item \textbf{An \code{Accept:\ application/json} request header makes
      the service double-encode its response.}  This was found while
      fixing the first two: the endpoint that still works for anonymous
      callers returned its payload as a JSON \emph{string} with every
      inner quote escaped, which the program's scanner cannot read.
\end{enumerate}

\noindent All three are addressed.  The program diagnoses its own
situation precisely, exits with a status that distinguishes ``sign in''
from ``try again later'', and retains a fully-working offline path.

\medskip
\noindent \textbf{Revision 5.0 of this report} adds
Section~\ref{sec:sibling}: the sibling \code{addswimuser} program is
broken, for a reason worth recording, and this program's roster now
reads from the database it writes to.

\medskip
\noindent \textbf{Revision 4.0} added
Section~\ref{sec:roster}: two defects found by tracing the path a new
swimmer takes into the database, both of which would have bitten the
first person to follow the instructions.

\medskip
\noindent \textbf{Revision 3.0} added
Section~\ref{sec:onedb}: the program now uses a single named database
and adapts to its schema rather than its own.

\medskip
\noindent \textbf{Revision 2.0} added
Section~\ref{sec:signin}, which closes the matter: rather than asking
an operator to extract a session from a browser, the program now
\emph{signs itself in} from a stored user name and password.  Doing so
required finding one more undocumented thing --- a newly-issued session
is refused by the times API until a separate ``permissions'' call has
been made against it --- and it exposed two further defects, both
introduced by the corrective action described above.  Live retrieval
works end to end.  Requirement-level verification stands at
\textbf{63 passed, 0 failed, 4 skipped}, against 38 passed and 21
skipped before this work.

% =================================================================
\section{Symptom and first observations}
% =================================================================

The reported symptom was total: no swimmer, no event, no output.

\begin{lstlisting}
$ ./swim-times -o stella
Error: person lookup request failed
$ echo $?
1
\end{lstlisting}

Three facts were established before forming any hypothesis.

\begin{enumerate}[leftmargin=*]
\item \textbf{The build was healthy.}  \code{ctangle} regenerated
      \code{swim-times.c} byte-for-byte identically to the committed
      copy, and \code{cc -O2 -Wall} compiled it without a warning.  The
      fault was not in the build.
\item \textbf{DNS and TLS were healthy.}
      \code{times-api.usaswimming.org} resolved (a Fastly CNAME) and
      completed a TLS handshake.  The fault was not connectivity.
\item \textbf{The service answered --- with a refusal.}  A hand-built
      \code{curl} request reproducing the program's headers returned
      \code{HTTP/2 403} with \code{content-length: 0}.  \emph{The
      program had received that 403 and called it a failed request.}
\end{enumerate}

\noindent That third observation split the investigation in two: why
does the service refuse, and why did the program not say so?

% =================================================================
\section{Root cause 1: the data hub closed its public endpoints}
% =================================================================

\subsection{Evidence}

The refusal is specific, not general.  Probing one endpoint from each
functional area, with identical headers, partitions the service's
surface cleanly:

\begin{center}
\begin{tabular}{llc}
\toprule
\textbf{Endpoint} & \textbf{Kind} & \textbf{Anonymous} \\
\midrule
\code{SearchFilter/GetAllEvents}            & reference & \textcolor{good}{200} \\
\code{SearchFilter/GetLscs}                 & reference & \textcolor{good}{200} \\
\code{TimesSearch/GetTimeStandardAgeGroups} & reference & \textcolor{good}{200} \\
\code{TimesSearch/GetMember/$\langle$id$\rangle$} & single member & \textcolor{good}{200} \\
\midrule
\code{TimesSearch/GetMembersForFilters}     & search & \textcolor{bad}{403} \\
\code{TimesSearch/BestTimes}                & times  & \textcolor{bad}{403} \\
\code{TimesSearch/GetBestTimesForMember}    & times  & \textcolor{bad}{403} \\
\code{TimesSearch/GetAllTimesForFilters}    & times  & \textcolor{bad}{403} \\
\code{TimesSearch/GetTopTimesLeaderBoard}   & times  & \textcolor{bad}{403} \\
\code{Ranking/Search}, \code{Records/SFSearch} & search & \textcolor{bad}{403} \\
\code{Meet/SFSearch}                        & search & \textcolor{bad}{403} \\
\bottomrule
\end{tabular}
\end{center}

Three further observations confirm this is deliberate policy rather
than an outage or a malformed request.

\paragraph{The \code{403} is keyed to the anonymous subject.}
Sending no \code{Usas-Sub-Id} at all yields \code{401 Unauthorized}.
Sending a fabricated subject also yields \code{401}.  Only the literal
\code{Anonymous} yields \code{403} --- the service recognises the
anonymous caller and \emph{then} declines.  A rejection for a malformed
body would be a \code{400}; a rejection for an unknown route would be a
\code{404}.  Both were observed from this service in other
circumstances, so it distinguishes them.

\paragraph{The front end says the same thing.}  Asking the data hub's
own authorisation service which routes an anonymous visitor may use:

\begin{lstlisting}
POST https://security-api.usaswimming.org/security/auth/GetDataHubSecurityInfoForIdp
     {"sub":"Anonymous","sid":""}
  -> { ..., "appRoutes": [ { "permissionType": "read", "routePath": "/" } ] }
\end{lstlisting}

\noindent A single route, \code{"/"}.  The browser application guards
\code{/search} with exactly this list and redirects an anonymous
visitor to the login page before it ever issues a times query.  The
API and the front end agree: the public individual-times search is
behind a login.

\paragraph{The lockout arrived in stages.}  The previous revision of
\code{swim-times.w} already records that
\code{GetAllTimesForFilters} and \code{GetSwimmerMeets} ``require a
signed-in subject and return \code{403} otherwise,'' while
\code{POST /BestTimes} was then still open.  \code{BestTimes} and
\code{GetMembersForFilters} have since joined them.  This is the third
narrowing of the same interface --- the Sisense JAQL API was retired
before it --- and it should be assumed to continue.

\subsection{Why this cannot be ``fixed''}

There is no remaining anonymous route to a swimmer's times.  Every
endpoint that returns a swim time is closed, and the one search that
would discover a \code{memberId} is closed with it.  No choice of
headers, body shape, \code{Device-Id}, \code{Origin}, or user agent
changes the outcome; each was tested individually.  A synthesised
\code{rate-key} does not help either, since the server verifies it
against a session it has issued.

The program can therefore only \emph{adapt}: identify the condition
precisely, offer a credentialed path for callers who have an account,
preserve the endpoints that still work, and degrade to the local
mirror.  Section~\ref{sec:fixes} describes what was done.

% =================================================================
\section{Root cause 2: the HTTP layer discarded the status code}
% =================================================================

This is the defect in \code{swim-times} proper, and the reason the
outage presented as a mystery rather than as a message.

\subsection{The defective code}

\begin{lstlisting}
static char *http_request(const char *url, const char *body)
{
    ...
    CURLcode rc = curl_easy_perform(curl);
    curl_slist_free_all(hdrs);
    curl_easy_cleanup(curl);

    if (rc != CURLE_OK) { free(buf.data); return NULL; }
    return buf.data;                      /* <-- NULL if body was empty */
}
\end{lstlisting}

\noindent and its caller:

\begin{lstlisting}
char *resp = http_request(url, body);
if (!resp) {
    fputs("Error: person lookup request failed\n", stderr);
    return NULL;
}
\end{lstlisting}

\subsection{The three faults in eleven lines}

\begin{enumerate}[leftmargin=*]
\item \textbf{The status code is never read.}  There is no call to
      \code{curl\_easy\_getinfo} anywhere in the original program.
      \code{curl\_easy\_perform} returns \code{CURLE\_OK} for a
      \emph{completed} transfer; a \code{403} is a completed transfer.
      The function's only notion of failure was therefore transport
      failure, and an HTTP-level refusal was invisible to it.
\item \textbf{\code{NULL} is overloaded.}  \code{buf.data} starts
      \code{NULL} and is allocated by the write callback on the first
      byte received.  A response with no body never fires the callback,
      so \code{buf.data} stays \code{NULL} and the function returns
      \code{NULL} --- the same value it uses to signal catastrophe.
      The data hub answers a refused request with
      \code{content-length: 0}.  The two cases were, quite literally,
      the same return value.
\item \textbf{The diagnostic asserts a cause it never established.}
      ``Person lookup request failed'' names the operation and nothing
      else.  It is true of a DNS failure, a TLS failure, a timeout, a
      \code{403}, a \code{500}, and an out-of-memory condition, and it
      distinguishes none of them.  A message that cannot be wrong is a
      message that cannot be useful.
\end{enumerate}

\noindent Two further weaknesses were found in the same function: no
\code{CURLOPT\_TIMEOUT} or \code{CURLOPT\_CONNECTTIMEOUT} (an
unresponsive host could stall a run indefinitely), and no redirect
following.

\subsection{The cost of the defect}

Had the status been reported, the original failure would have read
\code{HTTP 403 Forbidden} and the investigation would have begun at
the answer.  This is the general lesson: \emph{a diagnostic must report
what was observed, not what the author guessed.}  The observation ---
the status line --- was available for the cost of one function call and
was thrown away.

% =================================================================
\section{Root cause 3: \code{Accept} makes the service double-encode}
% =================================================================

While repairing the above, the \code{GetMember} path --- one of the few
still open to anonymous callers --- resolved the correct
\code{memberId} but reported the swimmer's name as \code{(unknown)}.
The response body was:

\begin{lstlisting}
"{\r\n  \"memberId\": \"6CD35348E5824C\",\r\n  \"fullName\": \"Katie ...
\end{lstlisting}

\noindent The whole document had been serialised a second time, as a
single JSON string with every inner quote escaped.  The scanner looks
for \code{"fullName"} and the bytes present are
\code{\textbackslash"fullName\textbackslash"}, so the match fails and
the field is silently skipped.

The trigger is a single request header.  Holding everything else
constant:

\begin{center}
\begin{tabular}{ll}
\toprule
\textbf{Headers sent} & \textbf{Response body} \\
\midrule
(base set)                              & plain JSON object \\
base $+$ \code{Content-Type: application/json} & plain JSON object \\
base $+$ \code{Accept: application/json}       & \textcolor{bad}{JSON string, escaped} \\
base $+$ both                                  & \textcolor{bad}{JSON string, escaped} \\
\bottomrule
\end{tabular}
\end{center}

\noindent The service's content negotiation treats an explicit
\code{Accept} as a request to serialise its already-serialised payload
again.  The data hub's own front end sends no \code{Accept} header, and
neither does the corrected program.  The rule this yields is worth
stating plainly: \textbf{send exactly the headers the service's own
client sends, and nothing more.}  A header added for tidiness changed
the response format.

% =================================================================
\section{Corrective actions}\label{sec:fixes}
% =================================================================

All changes are made in \code{swim-times.w}, the source of record;
\code{swim-times.c} is regenerated from it.

\subsection{The HTTP layer}

\begin{lstlisting}
static char *http_request(const char *url, const char *body, long *status)
{
    ...
    CURLcode rc = curl_easy_perform(curl);
    long code = 0;
    curl_easy_getinfo(curl, CURLINFO_RESPONSE_CODE, &code);
    ...
    if (rc != CURLE_OK) {
        fprintf(stderr, "Error: cannot reach %s: %s\n",
                url, curl_easy_strerror(rc));
        free(buf.data);
        return NULL;
    }
    if (status) *status = code;
    if (!buf.data) buf.data = calloc(1, 1);   /* "" is not NULL */
    return buf.data;
}
\end{lstlisting}

\noindent \code{NULL} now means one thing only: the request never
completed, and libcurl's own description of why has already been
printed.  Every response site checks \code{*status}.  Connect and total
timeouts are set; redirects are followed; content encodings are
negotiated; the \code{Accept} header is not sent.

\subsection{Accurate, actionable diagnostics}

\code{report\_http\_error} names the operation, the status, and the URL.
The first \code{401} or \code{403} of a run additionally prints the
remedy --- once, not once per event per swimmer:

\begin{lstlisting}
$ ./swim-times -o stella -e "100 FR SCY"
Error: member search failed --- HTTP 403 Forbidden
       https://times-api.usaswimming.org/swims/TimesSearch/GetMembersForFilters

  The USA Swimming Data Hub no longer answers this endpoint for
  anonymous callers.  Member search and every times query now
  require a signed-in account.  Two ways forward:

    1. Sign in at https://data.usaswimming.org, then copy the
       Usas-Sub-Id and Usas-Session-Id request headers the site
       sends to times-api.usaswimming.org and re-run with
           export USAS_SUB_ID=<subject>
           export USAS_SESSION_ID=<session>

    2. Re-run with  -o offline  to read times already mirrored in
       the local Postgres swim-times database.

  Run  -o diag  for a reachability table of the data-hub endpoints.
$ echo $?
77
\end{lstlisting}

\subsection{Meaningful exit codes}

\begin{center}
\begin{tabular}{rll}
\toprule
\textbf{Status} & \textbf{Name} & \textbf{Meaning} \\
\midrule
0  & \code{RC\_OK}      & success \\
1  & \code{RC\_ERROR}   & generic failure \\
2  & \code{RC\_USAGE}   & bad command line \\
69 & \code{RC\_UNAVAIL} & data hub unreachable or failed --- retry later \\
77 & \code{RC\_NOPERM}  & data hub refused for want of credentials --- sign in \\
\bottomrule
\end{tabular}
\end{center}

\noindent A wrapper script can now distinguish a transient outage from
a permissions problem without parsing English prose.

\subsection{Signed-in operation}

Credentials are read from the environment: \code{USAS\_SUB\_ID},
\code{USAS\_SESSION\_ID}, and an optional \code{USAS\_DEVICE\_ID}.
When a session id is present the program sends \code{Usas-Session-Id}
and a \code{rate-key} computed exactly as the data hub's front end
computes it:
\[
  \texttt{rate-key} \;=\;
  \left\lfloor \frac{t}{s} \right\rfloor \;\Vert\; \texttt{"."} \;\Vert\;
  \mathrm{HMAC\mbox{-}SHA256}_{\,\mathit{sid}}(t),
\]
where $t$ is the Unix time in milliseconds divided by $10^{4}$ and $s$
is the number of seconds elapsed since midnight UTC.

SHA-256 and HMAC-SHA256 are implemented in the literate source, in
about sixty lines, rather than taking on a TLS-library dependency for
one hash.  They are verified against the published vectors of
\code{FIPS 180-4} and \code{RFC 4231} by \code{-o selftest}, and the
suite cross-checks the expected digests against an independent
implementation (Python's \code{hashlib}).  Nothing is hard-coded: the
source contains no credential, and a test
(\req{REQ-D-11}) fails the build if one is introduced.

\subsection{A route around the closed search}

\code{GetMember/$\langle$id$\rangle$} remains open to anonymous
callers.  The \code{Swimmer} record therefore gained a sixth field,
\code{member\_id}; when it is set, the program confirms the swimmer
with that endpoint and never touches the closed search.  Katie
Ledecky's three roster rows carry her public identifier and use this
path today.  A new \code{-m $\langle$memberId$\rangle$} flag applies the
same short cut to any member, on or off the roster.

\subsection{Self-diagnosis}

\code{-o diag} reports the identity being presented and probes one
endpoint from each access class, so the three possible situations are
told apart at a glance:

\begin{lstlisting}
$ ./swim-times -o diag
swim-times data-hub diagnostics

  Usas-Sub-Id     : Anonymous
  Usas-Session-Id : (unset --- anonymous)
  Device-Id       : TGludXggLSBjdXJTGludsc3dpbSAtIDg2YWQ2YWFkYTg1MCAt...

  SearchFilter/GetAllEvents              GET  200 OK
  TimesSearch/GetTimeStandardAgeGroups   GET  200 OK
  TimesSearch/GetMember/<id>             GET  200 OK
  TimesSearch/GetMembersForFilters       POST 403 Forbidden
  TimesSearch/BestTimes                  POST 403 Forbidden

  200 = reachable; 401/403 = sign-in required (export USAS_SUB_ID and
  USAS_SESSION_ID).
\end{lstlisting}

\noindent All \code{200}: a working authorised client.  Reference
\code{200} with queries \code{403}: the anonymous lockout.  Nothing
answering: a genuine outage.  The test suite uses this same
classification to decide which tests can meaningfully run.

\subsection{Early abort instead of repeated failure}

\code{fetch\_pair} and \code{build\_member\_cache} now return a status,
and a refusal abandons the run.  Previously a locked-out full-roster
run would have issued roughly $29 \times 22 = 638$ doomed POSTs.  It now
issues one, reports once, and stops.

\subsection{Incidental repairs}

\begin{itemize}[leftmargin=*]
\item \textbf{Three \LaTeX{} errors eliminated.}  \code{pdftex} on the
      woven documentation reported three errors before this work (an
      underscore in a C comment, typeset in math mode) and reports none
      now.  \code{make doc} fails the build if any reappear.
\item \textbf{Module-size requirement restored.}  The \code{SWIMMERS}
      roster chunk was 38 lines, in violation of the standing
      twenty-four-line limit; it is split in two.  No chunk now exceeds
      the limit.
\item \textbf{Build no longer requires Docker.}  The \code{Makefile}
      uses a native \code{ctangle}/\code{cweave} when one is on
      \code{PATH} and falls back to the container wrapper otherwise.
      It also gained \code{doc}, \code{report}, \code{selftest},
      \code{diag}, \code{test}, and \code{db-test} targets.
\item \textbf{Schema drift recorded.}  The deployed \code{swim\_times}
      table predates \code{schema.sql}'s named constraint: its unique
      constraint is server-named and its timestamp column is
      \code{fetched\_at} rather than \code{loaded\_at}.  All seven
      columns the program writes are present and the \code{ON CONFLICT}
      tuple is correct, so behaviour is unaffected; \code{test-db.sh}
      now matches the constraint by column set rather than by name and
      reports the discrepancy instead of failing on it.
\end{itemize}

% =================================================================
\section{Verification}
% =================================================================

\begin{center}
\begin{tabular}{lrrr}
\toprule
\textbf{Suite} & \textbf{Pass} & \textbf{Fail} & \textbf{Skip} \\
\midrule
\code{test-swim-times.sh} & 71 & 0 & 4 \\
\code{test-db.sh}         & 20 & 0 & 0 \\
\bottomrule
\end{tabular}
\end{center}

\noindent The four skips are the refusal tests \req{REQ-E-01} through
\req{REQ-E-04}, which cannot run while the service is answering
normally.  Before the work of Section~\ref{sec:signin} the same host
recorded 38 passed and 21 skipped; the twenty-one were the live-data
tests, and signing in is what made them reachable.  Details and full
traceability are in the Software Test Plan, \code{CAST-STP-001}.

The tests that specifically pin the defects described here are:

\begin{center}
\begin{tabular}{lll}
\toprule
\textbf{ID} & \textbf{Guards against} & \textbf{Root cause} \\
\midrule
\req{REQ-H-01} & \code{http\_request} without a status out-parameter & 2 \\
\req{REQ-H-02} & status not read from libcurl                        & 2 \\
\req{REQ-H-03} & empty body conflated with failure                   & 2 \\
\req{REQ-H-04} & missing timeouts                                    & 2 \\
\req{REQ-H-05} & reintroduction of the \code{Accept} header          & 3 \\
\req{REQ-H-06} & an unchecked \code{http\_request} call site         & 2 \\
\req{REQ-E-01} & a refusal exiting with a generic status             & 1 \\
\req{REQ-E-02} & a diagnostic that does not name status and URL      & 1,\,2 \\
\req{REQ-E-03} & advice repeated once per event                      & 1 \\
\req{REQ-E-04} & a refusal replayed across the roster                & 1 \\
\req{REQ-S-01} & a wrong \code{rate-key} hash                        & --- \\
\req{REQ-G-03} & inability to classify the access tier               & 1 \\
\midrule
\req{REQ-L-03} & a credentials file being sourced rather than parsed  & --- \\
\req{REQ-L-04} & a foreign setting in that file being adopted         & --- \\
\req{REQ-L-06} & the activation step being dropped                    & --- \\
\req{REQ-L-07} & an anonymous fallback creeping back in               & --- \\
\req{REQ-F-34} & a \code{404} being treated as a failure              & --- \\
\req{DB-03}    & an override redirecting the database                & --- \\
\req{DB-08}    & a default appearing on \code{swim\_time\_id}        & --- \\
\req{DB-14}    & an elite label reaching the foreign key             & --- \\
\req{REQ-F-76} & an unstable surrogate id duplicating every swim     & --- \\
\req{REQ-F-77} & a stale roster entry aborting the whole walk        & --- \\
\req{DB-19}    & name spelling creating a duplicate person           & --- \\
\req{DB-20}    & a loose match merging two different swimmers        & --- \\
\req{REQ-R-02} & a database alias becoming unselectable again        & --- \\
\req{REQ-R-03} & a database row overriding a hand-tuned lookup       & --- \\
\bottomrule
\end{tabular}
\end{center}

% =================================================================
\section{Signing in}\label{sec:signin}
% =================================================================

Sections~2--6 left the program able to \emph{use} a session but not to
\emph{obtain} one: the operator was expected to sign in with a browser,
open the network inspector, and copy two request headers.  That is a
poor interface and a worse habit.  \code{requirements10.txt} asked for
the obvious alternative --- read a user name and password from
\code{/Users/evansjr/.usas-env}, sign in, and make that the default.
It is implemented, and it works.

\subsection{The flow}

Six requests, and then a seventh step that is the interesting part.

\begin{enumerate}[leftmargin=*]
\item \code{GET dhy-prod.usaswimming.org/bff/login} --- the data hub's
      back-end-for-front-end starts an OpenID~Connect
      authorization-code flow with PKCE and redirects to the identity
      provider's password form.
\item \code{POST login.usaswimming.org/Login} --- the form is
      resubmitted with \code{Username}, \code{Password}, the
      \code{ReturnUrl} it carried, and its antiforgery token.  Success
      is \code{302}; rejection is \code{200} with the form re-rendered.
\item \code{GET} the authorization callback.  Because the client
      registered \code{response\_mode=form\_post}, the reply is an HTML
      page whose body is a form a browser would submit by script.
\item \code{POST} that form's \code{code}, \code{state},
      \code{session\_state}, and \code{iss} to \code{/signin-oidc},
      which exchanges the code and sets the session cookie.
\item \code{GET /bff/userinfo} --- returns the claims, of which
      \code{sub} and \code{sid} are what the times API wants.
\item \code{POST security-api\dots/GetDataHubSecurityInfoForIdp} with
      that \code{sub} and \code{sid}.
\end{enumerate}

\subsection{The undocumented coupling}

Step 6 looks like a permissions query.  It returns the caller's name,
member id, and the list of application routes they may read, and
nothing about it suggests it changes anything.  It does.

\begin{center}
\begin{tabular}{lcc}
\toprule
\textbf{Request} & \textbf{Before step 6} & \textbf{After step 6} \\
\midrule
\code{GetMembersForFilters} & \textcolor{bad}{401} & \textcolor{good}{200} \\
\code{BestTimes}            & \textcolor{bad}{401} & \textcolor{good}{200} \\
\bottomrule
\end{tabular}
\end{center}

\noindent A session that has completed steps 1--5 and skipped step 6 is
answered \code{401 Unauthorized} by every protected times endpoint.
The same session, after one call to the security service, is answered
\code{200}.  This was established by running the two orders against
freshly-issued sessions repeatedly; it is not a timing effect, and
waiting does not substitute for it.

The effect is invisible from a browser, because the data hub's own
front end issues that call on start-up for its own reasons.  It cost
most of the time this section took to find, and it is worth recording
plainly: \textbf{obtaining a session is not the same as activating
one.}  \code{activate\_session} now performs step 6 on every run.

That diagnosis was possible only because of the work in
Section~\ref{sec:fixes}.  The distinction that mattered ---
\code{401} before, \code{403} when anonymous, \code{200} after --- is
precisely the distinction the original code discarded.

\subsection{What changed in the program}

\begin{itemize}[leftmargin=*]
\item \textbf{Signing in is the default and the only online path.}
      Every run that is not \code{-o offline} signs in before issuing a
      query.  There is no anonymous fallback: an unauthenticated search
      is a slower route to the same \code{403}, so the program does not
      attempt one.
\item \textbf{Credentials come from a file},
      \code{\$USAS\_ENV\_FILE} or \code{\$HOME/.usas-env}, in shell
      assignment syntax.  The file is \emph{parsed, never sourced} --- a
      command substitution in it is stored as text, not executed --- and
      only \code{USAS\_}-prefixed names are adopted.  A value already in
      the environment wins, so a single run can be overridden without
      editing the file.
\item \textbf{Failures name their cause.}  Absent credentials, rejected
      credentials, and a broken step each get their own message, and
      all exit \code{RC\_NOPERM}.  In particular a wrong password is
      reported as a wrong password, not as the \code{HTTP 200} the
      identity provider answers with.
\item \textbf{\code{-o diag} reports the sign-in} alongside the
      endpoint probes, and is the one mode that proceeds after a failed
      sign-in --- reporting what an unauthenticated caller can still
      reach is the diagnosis being requested.
\end{itemize}

\subsection{One deliberate omission}

The credentials file supplied also defines
\code{SWIM\_TIMES\_PGCONNINFO}, pointing at a database named
\code{swimming} on another host.  That database exists but holds a
different, normalised schema --- \code{swimmer}, \code{swim},
\code{meet}, \code{event} --- not the flat \code{swim\_times} table this
program writes.  Adopting the setting made \code{-o store} fail with
\code{relation "swim\_times" does not exist}.

The program therefore honours only \code{USAS\_}-prefixed names from
the file.  Silently retargeting this program's mirror to another
project's database is not a thing a credentials file should be able to
do by accident.  Exporting the variable still works, for an operator who
means it.

\subsection{Two defects this work exposed}

Both were introduced by the corrective action of
Section~\ref{sec:fixes} and both were found by the test suite on its
first authorised run --- which is the first time the suite had ever been
able to exercise these paths.

\paragraph{A \code{404} is data, not failure.}
\code{BestTimes} answers \code{404 Not Found} when the swimmer has
never swum that stroke at that distance.  For a ten-year-old that is
most of the thirty-one events --- the 1000 and 1650 free, the 200
butterfly, the 400 IM.  The abort-on-failure logic added in
Section~\ref{sec:fixes} treated it as fatal, so the default all-events
run ended at the first event the swimmer had never entered and printed
\emph{nothing}.  A \code{404} now contributes no rows and the walk
continues; \req{REQ-F-34} pins it and \req{REQ-F-21} is the end-to-end
check.

This is the mirror image of the original defect.  That one treated a
refusal as a transport failure; this one treated an empty result as a
refusal.  Both come from the same root: a status code carries
information, and each distinct status deserves a decision.

\paragraph{A chunk grew past the limit.}
The usage-text chunk reached twenty-six lines, over the standing
twenty-four-line rule, when the sign-in note was added.  It is split in
two.  \req{REQ-D-03} caught it.

\subsection{Cost}

Signing in adds six HTTP requests and roughly two seconds to every
online run.  No session is cached, because caching one would mean
writing a live session token to disk to save two seconds --- a poor
trade, and a staleness bug waiting to happen.  A single-event query
takes about three seconds end to end, well inside \req{REQ-P-02}.

% =================================================================
\section{One database, and not ours}\label{sec:onedb}
% =================================================================

Section~\ref{sec:signin} closed with a recommendation: keep the
Postgres mirror fed, because it is the only copy of the data that does
not depend on USA Swimming's access policy.  The follow-up instruction
made that concrete --- use the database on \code{100.77.243.69}, and
only that one.

\subsection{What was already there}

The database is not empty and this program did not create it.
\code{swimming} on that host holds 29 swimmers, 507 meets, and 4{,}482
swims, written by a sibling loader from USA Swimming's
\code{GetAllTimesForFilters} feed --- the richer endpoint that returns
every swim rather than only the bests.  Its schema is normalised: a
\code{swim} carries a \code{swimmer\_key} and a \code{meet\_key} rather
than repeating names, and its \code{event\_code},
\code{standard\_name}, and \code{age\_group\_label} columns are foreign
keys into reference tables.

That is not the flat \code{swim\_times} table this program had been
writing.  Two courses were open: create \code{swim\_times} alongside,
or adapt.  Adapting was chosen, and it is the better answer for a
reason worth stating --- a second table would have meant two copies of
the same swims in one database, and an \code{-o offline} that could see
only the copy this program wrote.  Adapting means \code{-o offline} now
reads the whole recorded history.  For Stella's 50~free short course
that is sixteen swims where the flat table had one.

\subsection{Three things the schema required}

\paragraph{Keys, not names.}
A swimmer must be resolved to a \code{swimmer\_key} before a swim can be
written.  The lookup matches on \code{member\_id} \emph{or}
\code{full\_name} because both are UNIQUE and three of the 29 existing
rows predate the member id: matching on only one of them finds nothing
and then collides on the other.  An existing row is refreshed with
\code{coalesce} rather than overwritten, so a fetch fills in a member
id that was missing without discarding anything already recorded.
Meets are resolved the same way, by name.

\paragraph{A primary key the feed does not supply.}
\code{swim.swim\_time\_id} normally holds USA Swimming's own
\code{swimTimeId}.  \code{BestTimes} does not return one.  The program
therefore derives a surrogate: the leading 63 bits of the SHA-256 of
the row's natural key, negated.  Negative values cannot collide with
the service's positive ids; the mapping is stable, so a re-fetch
reproduces the same id rather than inserting the swim again; and the
hash was already in the program for the \code{rate-key} of
Section~\ref{sec:signin}.  The sibling loader's own \code{BestTimes}
rows carry negative ids for the same reason.

\paragraph{A foreign key that does not know about elites.}
\code{time\_standard} holds the seven age-group levels --- \code{B}
through \code{AAAA} and \code{Slower Than B} --- and nothing else.
\code{BestTimes} also reports \code{Nats}, \code{Trials}, and
\code{Summer Jrs}, and writing one violates the foreign key.  Those map
to SQL \code{NULL}.  The label is not lost to the reader: it still
appears in the printed table and the CSV, neither of which is
constrained by anyone's schema.

\subsection{No override}

The connection string is a compiled-in constant and
\code{SWIM\_TIMES\_PGCONNINFO} is withdrawn.  This is deliberate and it
is the point of the instruction.  The database is shared; a run that
silently wrote somewhere else --- to a stale local copy, say, because a
variable was set in some other shell --- would look like it had worked
while contributing nothing.  There is now one place the data can go.

\subsection{Verification, and one deliberate intrusion}

Writing is idempotent, and that was demonstrated rather than asserted:
a store run against swims the sibling loader had already recorded left
the row count at 4{,}482 exactly, the bare \code{ON CONFLICT DO
NOTHING} absorbing every one.

That proves the guard but not the insert.  To show the insert actually
writes, one row --- Stella's 100~free, \code{swim\_time\_id}
339542386 --- was captured in full, deleted, re-fetched, and inspected.
The program resolved \code{swimmer\_key} 2 and the existing
\code{meet\_key} 13 rather than creating a duplicate meet, and wrote
the date, the seconds, the \code{BB} standard, the club, and
\code{source\_endpoint = 'BestTimes'}, \code{auth\_mode =
'authenticated'} under a negative surrogate id.  The original row was
then restored from the capture and the count verified back at 4{,}482
with its \code{power\_points} and \code{age\_group\_label} intact.

The test suites write only a clearly-marked fixture swimmer, removed by
an \code{EXIT} trap so an interrupted run still cleans up.  Everything
else they assert about the shared data is a read-only count.

\subsection{Two notes for the owner of that database}

\begin{enumerate}[leftmargin=*]
\item \textbf{Some swims are recorded twice already.}  The 20
      \code{BestTimes} rows the sibling loader wrote anonymously carry
      a \code{NULL} \code{swim\_date}, so the natural key does not match
      the dated row for the same swim from
      \code{GetAllTimesForFilters}.  Stella's 50~free appears both
      ways.  This predates the present work and nothing here has been
      changed to address it; deduplicating it is a decision for
      whoever owns the loader.
\item \textbf{This program writes a narrower row.}  \code{BestTimes}
      returns no \code{swimmerAge}, \code{powerPoints}, or
      \code{ageGroupLabel}, so rows it contributes leave those columns
      \code{NULL} where the richer feed fills them in.  Where both
      feeds cover a swim, the sibling loader's row is the better one ---
      which is exactly why this program never updates an existing row,
      only declines to duplicate it.
\end{enumerate}

% =================================================================
\section{Adding a swimmer, and two defects on the way}\label{sec:roster}
% =================================================================

The question ``how do I add a new swimmer to the database'' has a short
answer --- store a time for her --- and tracing it turned up two defects.
Both were found by inspection, not by a failing test, which is itself
worth recording: the suite had no case that walked a roster containing
an entry the service could not resolve, and none that wrote a swimmer
whose recorded name was spelled differently from the one the service
returns.

\subsection{There is no ``add swimmer'' operation}

A \code{swimmer} row appears only as a side effect of storing a swim:

\begin{lstlisting}
-o store -> db_insert_row() -> db_swimmer_key()
                                 |- db_find_swimmer()  SELECT ...
                                 |- found     -> UPDATE ... coalesce(...)
                                 `- not found -> INSERT ... RETURNING swimmer_key
\end{lstlisting}

\noindent A swimmer with no recorded times therefore never acquires a
row, and \code{-o fastest,store} records only the fastest time per
event.  Both follow from the design rather than contradicting it, but
neither is obvious from outside.

It is also worth noting that this program's roster is compiled in,
while the same database carries the sibling loader's roster as data, in
\code{swimmer\_alias} --- alias, search query, match substring, and age
window, the same five fields as the \code{SWIMMERS} table in
\code{swim-times.w}.  Neither reads the other, so the two drift.

\subsection{Defect: a \code{404} from the member search aborted the run}

\code{GetMembersForFilters} answers \code{404 Not Found} when no member
matches the name --- an ordinary outcome for a roster entry whose query
has gone stale.  \code{lookup\_member\_id} treated any non-\code{200} as
fatal and the dispatcher responded by breaking out of the roster loop.
The observed effect:

\begin{lstlisting}
$ ./swim-times -o stella,kalea,kenny,keith,bothwell-emma -e "50 FR SCY"
Swimmer: Stella Julianna Evans  (MemberId: 12E1458C85AF42)
Swimmer: Kalea Rose Benavente  (MemberId: 7667D74D6AD74D)
Error: member search failed --- HTTP 404 Not Found
\end{lstlisting}

\noindent Two swimmers reported, one stale entry, and the remaining two
silently dropped --- with an exit status suggesting the service had
failed.  Both \code{kenny} and \code{keith} had gone stale: the queries
\code{"Ray Evans"} and \code{"Santiago Evans"} match nothing any more.

This is the same mistake as the \code{BestTimes} \code{404} of
Section~\ref{sec:signin}, in the same shape: \textbf{a status meaning
``there is no such thing'' read as one meaning ``something went
wrong''.} Having fixed it once in the times path, we had not looked for
it in the search path.

The lookup now distinguishes the two outcomes with a flag, and the
dispatcher skips rather than breaks.  Skipped swimmers are counted and
the tally repeated at the end of the run, where it cannot be lost in
thirty screens of times.  A run in which nothing resolved exits
\code{RC\_UNAVAIL}; one that fetched some and skipped others exits
\code{RC\_OK} with the tally.  \code{kenny} and \code{keith} were
additionally given their member ids, which is the general remedy for a
stale query and cheaper besides --- one \code{GET} instead of a
\code{POST} and a scan.

\subsection{Defect: an exact name match would have duplicated people}

\code{db\_find\_swimmer} matched on \code{member\_id} or an exact
\code{full\_name}.  The two writers of this database learned their names
from different feeds: the sibling loader recorded middle
\emph{initials}, and the feed this program reads returns middle
\emph{names}.  Case differs too.

Checking all 24 resolvable roster swimmers against the database:

\begin{center}
\begin{tabular}{rl}
\toprule
\textbf{Swimmers} & \textbf{Outcome under the old match} \\
\midrule
22 & matched by \code{member\_id} --- correct \\
1  & matched by name; missing \code{member\_id} filled in --- correct \\
\textcolor{bad}{2} & \textcolor{bad}{no match $\rightarrow$ would insert a duplicate person} \\
\bottomrule
\end{tabular}
\end{center}

\noindent And the two were precisely the swimmers one would most want
to add, because they were the only roster members with \emph{zero}
recorded swims:

\begin{center}
\begin{tabular}{lll}
\toprule
\textbf{In the database} & \textbf{The service returns} \\
\midrule
\code{Liliya P Fedyshyn} (no member id)      & \code{Liliya Penny Fedyshyn} \\
\code{Addison F Wittmershaus} (no member id) & \code{Addison Marie Wittmershaus} \\
\bottomrule
\end{tabular}
\end{center}

The match now proceeds in three tiers --- member id, then full name
case-insensitively, then first and last name only --- with the second
and third confined to rows whose \code{member\_id} is \code{NULL}.
That guard is the substance of the fix, not a detail of it: a row
already carrying a different member id is a different person however
alike the names, and \textbf{merging two swimmers is a worse outcome
than duplicating one}.  \req{DB-20} asserts the guard directly, by
requiring that the same loose query match \emph{nothing} once the
fixture row carries a member id of its own.

A loose match links rather than renames.  \code{member\_id},
\code{lsc\_code}, and \code{club\_name} are filled in where absent;
\code{full\_name} is left as the database already had it, because a
loose match means the two spellings denote one person, not that ours is
the better spelling.

\subsection{Verification}

Storing for the two affected swimmers, which before the fix would have
created two duplicate people:

\begin{center}
\begin{tabular}{lrr}
\toprule
 & \textbf{Before} & \textbf{After} \\
\midrule
\code{swimmer} rows          & 29 & 29 \\
\code{swim} rows             & 4{,}482 & 4{,}496 \\
Roster members with no swims & 2 & 0 \\
\bottomrule
\end{tabular}
\end{center}

\noindent Both existing rows were linked to their member ids and given
their club, keeping the names the database already held.  Every one of
the 29 swimmers now has recorded times.

% =================================================================
\section{The sibling program, and one roster}\label{sec:sibling}
% =================================================================

Three sections of this report have referred to ``the sibling loader''
without naming it.  It is \code{addswimuser}, in
\code{\textasciitilde/myclaude2/addnewuser} --- itself a CWEB literate
program, with its own ConOps, SRS, SDD, test plan and test results.  It
admits a swimmer to the shared database:

\begin{lstlisting}
addswimuser -n <name of swimmer> -a <alias to use to look up swimmer>
\end{lstlisting}

\noindent writing a \code{swimmer} row and a \code{swimmer\_alias} row,
then populating that swimmer's times from
\code{GetAllTimesForFilters}.

\subsection{It no longer works}

\begin{lstlisting}
$ ./addswimuser -n "Maren Murch" -a mmurch-probe --dry-run
addswimuser: recipe: query "Murch", match "maren", club "College Area Swim Team"
warning: login step 3: no authorization code (bad username or password?)
warning: sign-in failed; continuing anonymously, which costs the age of every swim fetched
error: member search for "Murch": HTTP 403 (forbidden: subject not authorized)
\end{lstlisting}

\noindent The credentials are not wrong --- the same pair signs
\code{swim-times} in on the same machine.  Two defects, one of which is
fatal on its own.

\paragraph{The submit field was renamed.}
\code{addswimuser} posts \code{button=login} to the identity provider.
The form's field is now \code{loginButton}.  What makes this hard to
see is that the login endpoint \emph{accepts} the request either way
and answers \code{302}; the denial comes one redirect later, at the
authorization endpoint, as a \code{form\_post} page carrying
\code{error=access\_denied} instead of a code.  Holding everything else
constant:

\begin{center}
\begin{tabular}{ll}
\toprule
\textbf{Posted} & \textbf{Result after following the redirect} \\
\midrule
\code{button=login}      & \textcolor{bad}{\code{error=access\_denied}, no code} \\
\code{loginButton=login} & \textcolor{good}{code returned, login succeeds} \\
\bottomrule
\end{tabular}
\end{center}

\noindent Its own diagnostic --- ``no authorization code (bad username
or password?)'' --- is honest about what it saw and wrong about why,
which is the most expensive kind of error message.  The password was
never the problem.

\paragraph{And the session is never activated.}
\code{addswimuser} contains no call to
\code{GetDataHubSecurityInfoForIdp}.  Section~\ref{sec:signin} records
what that costs: a session that has completed the OpenID~Connect flow
is answered \code{401} by every protected times endpoint until that
call is made.  So fixing the field name alone would move the failure
from \code{403} to \code{401} without fixing it.

Both repairs are small and both are already written and proven in
\code{swim-times.w}; neither has been applied to \code{addswimuser},
which is a separate repository and was not in scope.

\subsection{One roster, in two places}

\code{addswimuser} writes \code{swimmer\_alias}, and that table turns
out to carry \emph{exactly} the fields this program's roster needs: the
alias, the \code{search\_query} sent to the data hub, the
\code{match\_substr} that picks the right candidate from the results,
and an optional age window.  Two programs, written independently, had
converged on the same five columns --- one holding them as a compiled-in
C table, the other as rows.

Neither read the other's, so they drifted.  \code{mmurch} has been in
the database since August with 141 recorded swims, and
\code{swim-times} could not select her.

The roster is now their union.  \code{SWIMMERS} is copied in first and
the database's aliases appended; a compiled-in entry wins any
collision, because its \code{match\_substr} has been tuned by hand and
behaviour should not change merely because a row was added elsewhere.
An alias that arrives from the database brings its \code{member\_id}
with it, which is the sturdiest lookup there is --- no search, and
nothing to go stale.

\begin{lstlisting}
$ ./swim-times -o list
swim-times roster: 33 aliases (29 compiled in, 4 from the database)

  alias                  swimmer                        memberId         source
  ---------------------- ------------------------------ ---------------- --------
  stella                 ~Julianna Evans / stella       (by search)      compiled
  ...
  cmadden                Clair Madden                   1970F79F814240   database
  ledecky                Katie Ledecky, all ages        6CD35348E5824C   database
  ledecky16              Katie Ledecky, age group 15-16 6CD35348E5824C   database
  mmurch                 Maren Murch                    29B0D16B40AD4D   database
\end{lstlisting}

\noindent \code{-o list} exists because the roster is no longer legible
by reading the source.  It needs the database but not the network, and
does not sign in.

The database is now opened on every run rather than only under
\code{store} and \code{offline}.  A database that cannot be reached is a
warning for every mode but \code{offline}, which has no other source of
times: refusing to fetch a swimmer this binary already knows about,
because a server at the far end of an overlay network is down, would be
a poor trade.

\subsection{Verification}

\code{mmurch} --- a database-only alias, never named in this program's
source --- fetches live, and \code{cmadden} reads offline, neither
requiring a rebuild.  \req{REQ-R-02} derives its expectation from
\code{swimmer\_alias} at run time rather than hard-coding a list, so it
is the case that would catch this regressing.

% =================================================================
\section{Operating the program now}
% =================================================================

\subsection{Set-up, once}

Put the account details in \code{\$HOME/.usas-env}:

\begin{lstlisting}
export USAS_USER="..."
export USAS_PASS="..."
\end{lstlisting}

\noindent \code{chmod 600} it.  Nothing else is required; the program
does the rest on every run.

\subsection{Ordinary use}

\begin{lstlisting}
./swim-times -o list                   # every selectable alias
./swim-times -o diag                   # expect: Sign-in ok, five 200s
./swim-times -o stella                 # all 31 events
./swim-times -o stella,kalea,csv,store # fetch, print CSV, and mirror
./swim-times -o kalea,fastest -e "100 FR SCY"
./swim-times -m 12E1458C85AF42         # straight to a memberId
\end{lstlisting}

\noindent A successful run announces whose account it used
(``\code{Signed in as John.}''), except in \code{csv} mode where a line
of prose would corrupt the output.

\subsection{Without credentials}

\begin{lstlisting}
./swim-times -o stella,offline         # read the shared database
\end{lstlisting}

\noindent \code{-o offline} needs no account and touches no network.  It
reads \code{swimming} on \code{100.77.243.69} --- the same database
\code{-o store} writes to, and the same one the sibling loader fills ---
so it reports the whole recorded history, not merely this program's own
contributions.  Keeping it fed remains the recommended routine: it is
the only copy of the data that does not depend on USA Swimming's
current access policy.

% =================================================================
\section{Risks and recommendations}
% =================================================================

\begin{enumerate}[leftmargin=*]
\item \textbf{Expect further narrowing.}  This is the third contraction
      of this interface.  \code{GetMember} being open today looks more
      like an oversight than a commitment.  \emph{Recommendation:} run
      \code{-o diag} before relying on a fetch, and treat the Postgres
      mirror as the system of record rather than as a cache.
\item \textbf{Session credentials are short-lived and personal.}  They
      are read from the environment and never written to disk by this
      program.  \emph{Recommendation:} keep them out of shell history
      and out of the repository; \req{REQ-D-11} fails the suite if a
      credential appears in the source.
\item \textbf{The \code{rate-key} path is now verified end to end.}
      A signed-in session was obtained and the assembled header
      accepted; all five \code{-o diag} probes return \code{200}.
\item \textbf{The password is stored in plaintext} in
      \code{\$HOME/.usas-env}, which is how the file was supplied and
      what \code{requirements10.txt} asked for.  The program reads it,
      never writes it, never logs it, and never sends it anywhere but
      the identity provider over TLS.  \emph{Recommendation:}
      \code{chmod 600} the file, and keep it out of the repository.
\item \textbf{The shared database has two writers.}  Neither deletes
      or updates the other's rows, and \code{source\_endpoint} plus
      \code{auth\_mode} tell them apart.  \emph{Recommendation:} keep
      it that way, and treat the sibling loader's
      \code{GetAllTimesForFilters} rows as authoritative where both
      feeds cover a swim.
\item \textbf{A full refresh is roughly 640 requests.}
      \emph{Recommendation:} run \code{./swim-times -o csv,store}
      across the roster on a schedule rather than ad hoc, so the
      database stays current without repeated manual passes.
\end{enumerate}

% =================================================================
\section{References}
% =================================================================

\begin{enumerate}[leftmargin=*,label={[\arabic*]}]
\item D.\,L.\,Parnas and P.\,C.\,Clements, ``A Rational Design Process:
      How and Why to Fake It,'' \textit{IEEE Trans.\ Software
      Engineering}, vol.\,SE-12, no.\,2, pp.\,251--257, Feb.\ 1986.
\item \textit{Secure Hash Standard (SHS)}, FIPS PUB 180-4, National
      Institute of Standards and Technology, Aug.\ 2015.
\item H.\,Krawczyk, M.\,Bellare, R.\,Canetti, ``HMAC: Keyed-Hashing for
      Message Authentication,'' RFC 2104, Feb.\ 1997.
\item M.\,Nystrom, ``Identifiers and Test Vectors for HMAC-SHA-224,
      HMAC-SHA-256, HMAC-SHA-384, and HMAC-SHA-512,'' RFC 4231,
      Dec.\ 2005.
\item R.\,Fielding, J.\,Reschke, eds., ``Hypertext Transfer Protocol
      (HTTP/1.1): Semantics and Content,'' RFC 7231, Jun.\ 2014
      (\S6.5.1 \code{400}, \S6.5.3 \code{403}, \S3.1.1.5
      \code{Content-Type}).
\item \textit{CAST Software Test Plan}, CAST-STP-001,
      \code{swim-times-stp.tex}.
\item \textit{CAST Software Requirements Specification}, CAST-SRS-001,
      \code{swim-times-srs.tex}.
\item \code{swim-times.w} --- the literate source of record, whose
      woven form (\code{swim-times.pdf}) documents every change
      described here at the point it occurs.
\end{enumerate}

\end{document}
